% Replacement passages for the supplied manuscript; the original PDF is unchanged.

% Abstract/conclusion: distinguish horizon dependence from full minimax optimality.
We derive a high-probability $\widetilde O(K^{-1/2})$ guarantee after a sufficient
burn-in period, matching the horizon dependence of a worst-case expected-risk
lower bound. The bounds retain different dependence on other problem parameters.

% Section 4, comparison of the displayed upper and lower bounds.
The upper and lower bounds both have inverse-square-root dependence on the
number of online interactions, up to logarithmic factors in the upper bound.
The upper bound retains additional dependence on dimension, inverse temperature,
the parameter radius, and confidence. The displayed expected-risk lower bound
does not establish optimal dependence on the confidence parameter.

% Section 5, scope of empirical evidence.
We evaluate the practical algorithm in a GUARD-derived anti-poaching game with
a correctly specified simulated logit follower. The experiments compare complete
data-collection and estimation protocols: active methods choose leader policies,
fixed-expert MLE receives observations under a source-optimal leader, and the
OLS baselines use a joint-QRE expert. Zero-sum OLS also observes leader actions.
The source payoff is used by the evaluator and demonstration oracle, not as an
input to the reward estimator. If both that payoff and its exact zero-sum
relationship were available to the learner, no source reward estimation would
be necessary.
The results concern this game and the evaluated downstream objectives; they do
not establish superiority across arbitrary games or robustness to model
misspecification.

% Appendix I.2 and Figure 8. Preserve the existing heatmap.
Figure~\ref{fig:appendix-lobeke-hyperparameters} reports post-selection sensitivity
to the warm-up budget and online information regularizer. The grid crosses
warm-up budgets $\{0,50,200\}$ with multipliers
$\lambda/\eta^2\in\{0.003,0.01,0.03,0.3,1\}$, using $1{,}600$ follower observations,
eight paired learning seeds, and $100$ fixed normalized transfer tasks. MLE ridge
is zero and is distinct from the online information regularizer. Outlines mark
the retained protocols, not the lowest cells in this later sensitivity grid.
They were selected in an earlier ridge-$10^{-4}$ stage and were not reselected
from this figure. The subsequent ridge-zero replay reused an earlier confirmation
split. These analyses do not establish a unique optimal regularizer.
The retained multipliers refer to the numerical inverse temperature
$\eta\approx2.165$, giving $\lambda=0.0140625$ with warm-up and $0.140625$ without.

% Appendix I.3.1: complete sparse-information setup and interpretation.
Let $\mathcal A=[m]$, $\mathcal B=[2]$, and
$r^2_\omega(a,b)=\omega\mathbf{1}\{a=1,b=2\}$, with true coefficient $\omega=1$
and known inverse temperature $\eta=2$. The leader earns
$r^1(1,b)=\mathbf{1}\{b=2\}$ and $r^1(a,b)=3/4$ for $a>1$.
Only the first pure leader action is informative. A mixed policy assigning it
mass $x$ induces response probability $p=\sigma(\eta\omega x)$ and Fisher information
\begin{equation*}
I_x(\omega)=\eta^2x^2p(1-p).
\end{equation*}
Information is zero exactly when $x=0$; random mixed policies generally provide
positive but diluted information. At the initial estimate $\hat\omega=0$, the
ratio of maximal information to information averaged over a uniform-simplex
policy is $m(m+1)/2$. The online information optimizer can choose an interior
mass at other estimates, so warm-up and online selection need not always play
the same pure action. We report the no-warm-up variant to isolate online selection.

% Appendix I.3.2 and Figure 10: interpretation of the hard-family experiment.
The hard family changes with the observation budget:
$\omega=\epsilon e_i$, where $\epsilon=\tfrac12\sqrt{d/K}$.
A uniform leader policy that ignores the observations already has loss
\begin{equation*}
\operatorname{SubOpt}(\mu_{\mathrm{uniform}})=\sigma(\epsilon)-\sigma(\epsilon/d)=O(\sqrt{d/K}).
\end{equation*}
The empirical curves therefore illustrate behavior on this shrinking-separation
family. Their slope alone does not demonstrate learning on a fixed instance,
prove a lower bound over all algorithms, or establish minimax optimality.
The theoretical guarantee rests on the proof, independently of this illustration.
The experiment uses $d\in\{33,64\}$ and $16$ independent replicates of every active
coordinate; the uniform-coordinate mean therefore includes $16d$ trajectories
per budget and dimension.

% Model description: retain the normalization-induced temperature rescaling.
We use $\phi=\phi_{\mathrm{raw}}/L$, where
$L=\max_{a,b}\lVert\phi_{\mathrm{raw}}(a,b)\rVert_2$, together with
$\eta=L\eta_{\mathrm{raw}}$ and $\eta_{\mathrm{raw}}=1.25$.
This preserves response probabilities while keeping the reward coefficients
$\omega^\star=(0,1,-2/3)$ unchanged. The numerical inverse temperature is
approximately $2.165$ in this instance.

% Optional implementation justification for the hard-family vertex reduction.
On the hard family, the information objective admits a vertex maximizer when
$\lVert\hat\omega\rVert_2\leq1$ and $\eta=1$. Let
$v(z)=\sigma(z)(1-\sigma(z))$ and let $G$ be the diagonal regularized information
matrix produced by previous vertex observations. The objective is
$v(\mu^\top\hat\omega)\sum_i\mu_i^2/G_{ii}$.
On $[-1,1]$, $|\frac{d}{dz}\log v(z)|\leq\tanh(1/2)$, while
$|\mu^\top\hat\omega-\hat\omega_i|\leq2(1-\mu_i)$. For $\mu_i>0$,
\begin{equation*}
\log\frac{\mu_i v(\mu^\top\hat\omega)}{v(\hat\omega_i)}\leq\log\mu_i+2\tanh(1/2)(1-\mu_i)\leq0.
\end{equation*}
The inequality also holds without logarithms at $\mu_i=0$. Consequently,
\begin{equation*}
v(\mu^\top\hat\omega)\sum_i\frac{\mu_i^2}{G_{ii}}\leq\sum_i\mu_i\frac{v(\hat\omega_i)}{G_{ii}}\leq\max_i\frac{v(\hat\omega_i)}{G_{ii}}.
\end{equation*}
Evaluating all vertices therefore gives a global maximizer for this specialized
instance; this does not certify global policy optimization in the Lob\'ek\'e game.

% Additional fixed-settings confirmation, separate from the retained figures.
We additionally evaluated the two retained Active protocols on $50$ fresh
learning seeds and $100$ fixed new transfer tasks, using $1{,}600$ follower
observations per run. Mean transfer SubOpt was $0.000286$ without warm-up and
$0.000261$ with warm-up, compared with $0.003488$ for ridge-zero fixed-expert MLE.
The paired differences were $-0.003202$ and $-0.003227$, with simultaneous
$95\%$ intervals $[-0.004272,-0.002133]$ and $[-0.004277,-0.002177]$ for the two
comparisons. These intervals are over learning seeds, conditional on the task
suite. The full multistart deployment evaluation repairs missed solutions from
the initial fast numerical evaluation. Matched-regularizer warm-up comparisons
were inconclusive. An exploratory expert-ridge sweep substantially reduced the
baseline loss, motivating a separate fresh confirmation against its selected
setting; the large ridge-zero comparison should not be presented as an advantage
over all regularized expert estimators.

% Independent selected-baseline check. Do not claim adjusted significance.
We selected fixed-expert MLE ridge $10$ from $\{0,0.01,0.1,1,10\}$ on the first
confirmation suite, then evaluated it against the unchanged Active protocols
on another $50$ learning seeds and $50$ transfer tasks. Mean SubOpt was
$0.000251$ for no-warm-up Active, $0.000205$ for warm-up Active, and $0.000347$
for the selected expert baseline. The paired differences were $-0.000096$
and $-0.000142$, with simultaneous $95\%$ intervals
$[-0.000262,0.000069]$ and $[-0.000286,0.000003]$ for the two comparisons.
Holm-adjusted $p$-values were $0.185$ and $0.055$. Thus regularization closes
much of the unregularized baseline gap, and this independent check does not
establish superiority over the selected regularized expert. The inference is
conditional on these synthetic tasks in one fixed game.
